\subsection{Household-level consequences of overreaction}%
% duplicate label removed (stale file): sec:chfs_consequences

The welfare cost calculation in Section~\ref{sec:policy} is
model-calibrated: I feed the estimated overreaction parameter into a
three-equation New Keynesian model and compute the implied loss. This
approach disciplines the magnitude but says nothing about whether
individual households who overreact actually behave differently. The
CHFS panel provides direct, if limited, evidence.

\paragraph{Classifying overreactors.}
I residualize the 2011 A4008 inflation-expectation measure on province
fixed effects and controls (log income, urban indicator). The residual
captures the household-specific component of inflation expectations
that is not explained by location or observables. I define
$Overreactor_{i,2011}$ as an indicator for households in the top
quartile of this residual distribution. These are households whose
inflation expectations in 2011 were substantially higher than what
their province and demographics would predict, the operational
signature of diagnostic overweighting of salient price signals.

\paragraph{Forward-tracking specification.}
I link the 2011 belief classification to income changes observed in
subsequent waves:
\begin{equation*}% duplicate label removed (stale file): eq:chfs_income
  \Delta \log Y_{i,t} = \alpha + \beta \cdot Overreactor_{i,2011}
    + \gamma X_{i,2011} + \delta_p + \varepsilon_i,
\end{equation*}
where $\Delta \log Y_{i,t}$ is the change in log total household
income between 2011 and wave $t \in \{2013, 2015\}$, $X_{i,2011}$
includes log income and an urban indicator, and $\delta_p$ denotes
province fixed effects. Consumption data are not consistently
available across CHFS waves, so I use income growth as the outcome.
The prediction follows from standard consumption-savings theory under
diagnostic beliefs: households who overestimate inflation in 2011
may front-load consumption or misallocate savings, and these
suboptimal decisions could depress subsequent income growth through
reduced human-capital or business investment.

\paragraph{Results.}
Table~\ref{tab:chfs_consumption_consequences} reports the estimates
at two horizons. At the 2011--2013 horizon, $\hat{\beta} = 0.022$
(SE $= 0.051$, $N = 6{,}640$), small and statistically
indistinguishable from zero. At the 2011--2015 horizon,
$\hat{\beta} = 0.106$ (SE $= 0.060$, $N = 5{,}550$), marginally
significant at the 10\% level. The positive sign at longer horizons
is unexpected under the simple front-loading prediction, which implies
a negative coefficient. One interpretation is that overreactors
in 2011 were disproportionately drawn from households experiencing
genuine local price pressure, and the residualization on province
fixed effects and observables does not fully absorb this. Higher
inflation expectations in 2011 may proxy partly for local economic
dynamism that also predicts faster income growth through 2015.

\begin{table}[htbp]
\centering
\begin{threeparttable}
\caption{Behavioral Consequences of Diagnostic Overreaction: Forward Income}
\label{tab:chfs_consumption_consequences}
{\small
\begin{tabular}{lcc}
\toprule
 & (1) & (2) \\
 & 2011--2013 & 2011--2015 \\
\midrule
\multicolumn{3}{l}{\textit{Panel A: Dependent variable = change in log household income relative to 2011}} \\
\addlinespace
Overreactor & $0.022$ & $0.106^{*}$ \\
 & $(0.051)$ & $(0.060)$ \\
 & $[0.666]$ & $[0.077]$ \\
\addlinespace
Observations & 6,640 & 5,550 \\
Province FE & Yes & Yes \\
Controls & Yes & Yes \\
\bottomrule
\end{tabular}
}
\begin{tablenotes}[flushleft]
\footnotesize
\item \textit{Notes:} The dependent variable is the change in log household income relative to 2011. Overreactor is defined as the top quartile of province-residualized inflation expectation in 2011. Controls include log income and an urban indicator. Province fixed effects are included. Heteroskedasticity-robust standard errors are in parentheses; p-values are in brackets. $^{*}p<0.10$; $^{**}p<0.05$; $^{***}p<0.01$.
\end{tablenotes}
\end{threeparttable}
\end{table}


\paragraph{Interpretation and limitations.}
The forward-tracking evidence is suggestive rather than dispositive.
The null result at the two-year horizon and the marginally significant
positive coefficient at the four-year horizon do not confirm the
simple Euler-equation prediction that overreactors experience
consumption or income declines. Several factors limit statistical
power. The CHFS panel uses income rather than consumption as
the outcome, and income is a noisier proxy for the welfare channel.
The overreactor classification relies on a single
cross-sectional wave with a coarse 5-point expectation scale, which
attenuates the signal. Sample attrition between 2011 and 2015
also reduces the estimation sample from 6,640 to 5,550 households,
potentially introducing selection.

\paragraph{Connection to the NK welfare calculation.}
The model-calibrated welfare loss in Section~\ref{sec:policy} implies
that diagnostic overreaction at $\theta = 0.55$ raises the consumption
variance component of welfare loss by 37.6\% relative to rational
expectations. The household-level evidence here does not contradict
that calibration, but it does not independently confirm it either.
The forward-tracking test lacks statistical power to detect the
income effects that the model predicts. The model provides the
aggregation logic through general-equilibrium amplification; the
household data confirm that overreaction is a measurable
cross-sectional phenomenon (Section~\ref{sec:chfs_info_channel})
even if its forward consequences are imprecisely estimated with the
available panel structure.




